\section{Implementasi Spesifikasi Bonus}\label{sub:Implementasi Spesifikasi Bonus}

\subsection{Algoritma Pathfinding Alternatif}\label{sub:Algoritma Pathfinding Alternatif}

Bonus algoritma pathfinding alternatif dipenuhi dengan menambahkan BFS
dan DFS. Keduanya dapat dijalankan dari CLI maupun GUI. Pada GUI,
frontend mengirim nilai \texttt{BFS} atau \texttt{DFS} melalui field
\texttt{algorithm}, lalu backend memanggil \texttt{BFSWithTrace} atau
\texttt{DFSWithTrace}.

\begin{lstlisting}[language=Go, basicstyle=\ttfamily\small, breaklines=true]
case "BFS":
    result, trace := solver.BFSWithTrace(board)
    return result, trace, nil
case "DFS":
    result, trace := solver.DFSWithTrace(board)
    return result, trace, nil
\end{lstlisting}

Kedua algoritma ini tidak memerlukan heuristic, sehingga dropdown
heuristic pada GUI hanya aktif untuk GBFS dan A*.

\subsection{GUI Web}\label{sub:GUI}

GUI dibuat sebagai aplikasi web agar solver dapat digunakan tanpa
interaksi terminal. Arsitektur GUI terdiri dari dua bagian.
\begin{enumerate}
    \item Backend Go pada \texttt{src/web/backend/server.go}.
    \item Frontend React pada \texttt{src/web/frontend}.
\end{enumerate}

Backend menyediakan endpoint \texttt{POST /solve}. Endpoint ini
menerima \texttt{multipart/form-data} berisi file testcase, nama
algoritma, dan mode heuristic. Setelah request diterima, backend
menjalankan parser dan solver, lalu mengembalikan respons JSON.

\begin{lstlisting}[language=bash, basicstyle=\ttfamily\small, breaklines=true]
POST /solve
file=<testcase.txt>
algorithm=UCS|GBFS|ASTAR|BFS|DFS
heuristic=H1|H2|H3
\end{lstlisting}

Respons backend berisi informasi berikut.
\begin{itemize}[label=\textbullet]
    \item \texttt{algorithm}: algoritma yang dijalankan.
    \item \texttt{heuristic}: heuristic yang dipakai untuk GBFS atau A*.
    \item \texttt{found}: status apakah solusi ditemukan.
    \item \texttt{moves}: urutan gerakan solusi.
    \item \texttt{totalCost}: total biaya solusi.
    \item \texttt{iterations}: jumlah simpul yang diekspansi.
    \item \texttt{executionMs}: waktu eksekusi solver dalam milidetik.
    \item \texttt{steps}: papan visual untuk setiap langkah solusi.
    \item \texttt{trace}: node dan edge untuk visualisasi graf pencarian.
\end{itemize}

\subsection{Frontend React}\label{sub:Frontend React}

Frontend menyediakan beberapa komponen utama.
\begin{enumerate}
    \item \texttt{ControlPanel.jsx}: form upload file, pilihan
        algoritma, dan pilihan heuristic.
    \item \texttt{Board.jsx}: visualisasi grid puzzle.
    \item \texttt{ResultPanel.jsx}: ringkasan hasil solver.
    \item \texttt{PlaybackControls.jsx}: kontrol previous, next,
        play/pause, slider step, dan speed.
    \item \texttt{GraphView.jsx}: visualisasi trace pencarian dalam
        bentuk graf.
\end{enumerate}

Alur penggunaan GUI adalah sebagai berikut.
\begin{enumerate}
    \item Pengguna mengunggah file testcase \texttt{.txt}.
    \item Pengguna memilih algoritma UCS, GBFS, A*, BFS, atau DFS.
    \item Jika memilih GBFS atau A*, pengguna memilih heuristic.
    \item Frontend mengirim request ke endpoint \texttt{/solve}.
    \item Backend mengembalikan hasil solver.
    \item Frontend menampilkan hasil, playback solusi, graf
        pencarian, dan tombol \texttt{Save .txt} untuk menyimpan
        ringkasan solusi.
\end{enumerate}

\subsection{Trace Graf Pencarian}\label{sub:Trace Graf Pencarian}

Untuk kebutuhan GUI, solver memiliki varian fungsi \texttt{WithTrace}.
Trace tidak dimasukkan ke dalam \texttt{solver.Result} agar struktur
hasil utama tetap fokus pada solusi pencarian. Data trace diletakkan
pada package terpisah, yaitu \texttt{searchtrace}.

\begin{lstlisting}[language=Go, basicstyle=\ttfamily\small, breaklines=true]
type Step struct {
    Expanded Node
    Edges    []Edge
}
\end{lstlisting}

Setiap \texttt{Step} merepresentasikan satu ekspansi simpul. Edge pada
trace menyimpan informasi arah, biaya langkah, nilai \texttt{g},
nilai \texttt{h}, prioritas, dan apakah edge tersebut diterima masuk
frontier. Dengan informasi ini, GUI dapat membedakan successor yang
dipakai solver dan successor yang ditolak karena sudah ada rute lebih
murah.

\subsection{Alternatif Heuristic}\label{sub:Alternatif Heuristic}

Bonus alternatif heuristic dipenuhi dengan menyediakan \texttt{H2} dan
\texttt{H3} selain heuristic utama \texttt{H1}. Ketiga heuristic dapat
dipilih pada GUI ketika algoritma yang digunakan adalah GBFS atau A*.

\subsection{Docker}\label{sub:Docker}

Dockerfile menggunakan multi-stage build.
\begin{enumerate}
    \item Stage Node membangun frontend React.
    \item Stage Go membangun binary backend.
    \item Stage Alpine menjalankan backend dan menyajikan hasil build
        frontend sebagai static SPA.
\end{enumerate}

Konfigurasi \texttt{docker-compose.yml} di root repository mengatur
build image dari \texttt{src/Dockerfile} dan memetakan port
\texttt{8080}. Dengan desain ini, pengguna cukup menjalankan
\texttt{docker compose down} lalu \texttt{docker compose up --build}
untuk mengakses GUI lengkap melalui browser.
